\section{Silogisme Hipotetis dan Silogisme Disjungtif}

Keluarga penarikan argumen Silogisme mewujud bagai jembatan titian logika. Jika Duo Modus (Ponens dan Tollens) berinteraksi dengan pernyataan absolut berstatus pasti, klan Silogisme bermain cantik mentransmisikan nilai logika antar beberapa persimpangan bersyarat maupun pilihan percabangan \cite{rosen2019}.

\subsection{Silogisme Hipotetis (Hukum Rantai Implikasi)}

Bentuk inferensi pengalir ini menculik Sifat Transitif murni yang banyak diemban fungsi relasional matematika. 

\begin{teorema}[Silogisme Hipotetis]
Jika fenomena $p$ digaransi menyulut meletusnya rantai fenomena $q$. Lantas, fenomena penyerta $q$ tersebut secara inheren digaransi bakal merembes mengakibatkan rentetan fenomena akhir $r$. Maka rasio menyimpulkan secara ajaib jalan pintas komando (bypass): $p$ sanggup memicu terjadinya konklusi akhir $r$ secara tidak langsung (Hipotetis berantai).
\end{teorema}

\textbf{Anatomi Argumen Silogisme Hipotetis:}
\[
\begin{array}{c}
p \rightarrow q \\
q \rightarrow r \\
\hline
\therefore p \rightarrow r
\end{array}
\]
Perhatikan elemen mediator penengah $q$ dieliminasi dari sorotan panggung kesimpulan sembari ia mewariskan pilar energi transitif dari $p$ meluncur tembus ke leher $r$.

\begin{contoh}
Argumen rantai motivasi belajar:
\begin{itemize}
    \item \textbf{Premis 1 ($p \rightarrow q$)}: "Bila saya mengerahkan tekad merajut buku Logika di malam hari, maka gembok materi aljabar proposisi bakal berhasil dicerna otak."
    \item \textbf{Premis 2 ($q \rightarrow r$)}: "Bila gembok aljabar proposisi lumat tercerna otak, maka besok fajar saya pasti lulus cemerlang mata kuliah bapak dosen."
    \item \textbf{Kesimpulan ($\therefore p \rightarrow r$)}: "Jadi, bila saya mengerahkan tekad merajut buku Logika malam ini, maka besok fajar saya sanggup merengkuh lulus cemerlang."
\end{itemize}
\end{contoh}

\subsection{Silogisme Disjungtif (Pemangkasan Cabang Alternatif)}

Ketika kita tertohok di ujung jalan persimpangan Y di mana salah satu cabang menjanjikan tebing curam kematian, manusia rasional hanya punya satu sisa nasib tersisa.

\begin{teorema}[Silogisme Disjungtif]
Berangkat dari ketegasan fondasi Disjungsi Or ($p \lor q$)—yang memastikan sekurang-kurangnya sebuah batang klausa berhak mengemban mahkota Benar T. Di kala nasib berputar dan membongkar bahwa cabang rute $p$ mendadak cacat roboh tak bernilai ($\neg p$), praktis alam semesta menitipkan keseluruhan nasib nilai T murni tanpa sanggahan dipanggul sisa cabang rute sendirian: Opsi $q$ \textbf{DIWAJIBKAN MUTLAK} untuk menjelma Benar!
\end{teorema}

\textbf{Anatomi Argumen Silogisme Disjungtif Asimetris:}
\[
\begin{array}{ccc}
p \vee q & & p \vee q\\
\neg p   & \text{Atau varian sisi berlawanannya:} & \neg q \\
\hline
\therefore q & & \therefore p
\end{array}
\]

\begin{contoh}
Skenario kepastian alibi kriminalisasi:
\begin{itemize}
    \item \textbf{Premis 1 ($p \vee q$)}: "Oknum maling misterius melarikan diri menjebol jendela ($p$), ATAU maling tersebut lari menyelinap melalui lubang plafon atap ($q$)."
    \item \textbf{Premis 2 ($\neg p$)}: "Investigasi di TKP menyatakan kusen maupun teralis seluruh jendela terkunci utuh segel tak ternoda (Maling \textbf{mustahil} lari pecahkan jendela)."
    \item \textbf{Kesimpulan ($\therefore q$)}: "Fiks valid tiada tawar menawar! Maling itu sungguh menyelinap lari menggasak plafon atap."
\end{itemize}
\end{contoh}
Bentuk argumen raksasa dari argumen ala detektif Conan ini: $[(p \lor q) \land \neg p] \rightarrow q$. Silakan bedah melalui iterasi Tabel Kebenaran 4 baris, niscaya keagungan gelar Tautologi terpatri kokoh sepadan dengannya.
